% Part: normal-modal-logic
% Chapter: sequent-calculus
% Section: introduction
%READERNOTE{SOL6-A183: S5 साठी कट-मुक्त कलन ज्ञात नसल्याचा मूळ दावा ऐतिहासिक मसुद्यापुरता केला. प्रत्यक्ष वाचलेल्या Aghaei आणि Mohammadi यांच्या 2018 लेखाचा स्रोत जोडला; त्यातील G3S5 आणि पुढे दिलेला विशिष्ट LK-विस्तार यांची व्याप्ती वेगळी ठेवली.}

\documentclass[../../../include/open-logic-section]{subfiles}

\begin{document}

\olfileid{nml}{seq}{int}

\olsection{प्रस्तावना}

विधानीय तर्कशास्त्राच्या क्रमवर्ती कलनात
~\iftag{prvBox}{$\Box$\iftag{prvDiamond}{ आणि~}{}}{}%
\iftag{prvDiamond}{$\Diamond$}{} यांच्याशी संबंधित
अतिरिक्त नियम घालून त्याचा विस्तार करता येतो.
उदाहरणार्थ,~$\Log{K}$ साठी \Log{LK} बरोबर पुढील
नियम मिळतात:
\[
  \iftag{prvBox}{\iftag{prvDiamond}{% <> and [] primitive
    \Axiom$\Gamma \fCenter \Delta, !A$
    \RightLabel{$\Box$}
    \UnaryInf$\Box\Gamma \fCenter \Diamond\Delta, \Box!A$
    \DisplayProof
    \qquad
    \Axiom$!A, \Gamma \fCenter \Delta$
    \RightLabel{$\Diamond$}
    \UnaryInf$\Diamond!A, \Box\Gamma \fCenter \Diamond\Delta$
    \DisplayProof
}{% only Box primitive
\Axiom$\Gamma \fCenter !A$
\RightLabel{$\Box$}
\UnaryInf$\Box\Gamma \fCenter \Box!A$
\DisplayProof
}}{% only <> primitive
  \Axiom$!A \fCenter \Delta$
  \RightLabel{$\Diamond$}
  \UnaryInf$\Diamond!A \fCenter \Diamond\Delta$
  \DisplayProof
}
\]
~$\Log{K}$ च्या विस्तारांसाठीही आणखी नियम
जोडावे लागतात.

प्रत्येक मोडल तर्कशास्त्रासाठी अशा प्रकारचे
क्रमवर्ती कलन उपलब्ध असेलच असे नाही. अर्थदृष्ट्या
साध्या असलेल्या~$\Log{S5}$ साठीसुद्धा (त्याची
व्याख्या प्राप्यता-संबंध मुळीच न वापरता करता येते),
मूळ मसुद्यात~$\Log{LK}$ पासून मिळणारे आणि~\Cut
नियमाशिवाय संपूर्ण असणारे कलन ज्ञात नसल्याचे म्हटले आहे.
हे सध्याच्या सर्व क्रमवर्ती कलनांविषयीचे विधान समजू नये.
उदाहरणार्थ, \href{https://arxiv.org/abs/1711.04634v3}{Aghaei
आणि Mohammadi यांच्या 2018 मधील लेखात} $\Log{G3c}$ चा
विस्तार असलेले $\Log{G3S5}$ दिले आहे. कलम~3 मधील नियम
आणि प्रमेये~5.1 व~5.2 त्यातील कटची अनुमतता आणि संपूर्णता
देतात; सिद्धतेसाठी त्याच कलनाचे अर्धविराम वापरणारे रूप घेतले आहे.
या प्रकरणात पुढे दिलेल्या विशिष्ट $\Log{LK}$-विस्तारासाठी
कटची गरज असल्याचा निष्कर्ष वेगळा आहे.
$\Log{S5}$ साठी कट-मुक्त संपूर्ण \emph{हायपरक्रमवर्ती}
कलनही आहे.

\end{document}
